\section{Related Work}

Our work builds on three research threads: weight-space learning for property prediction, permutation-equivariant neural architectures, and inductive biases in deep learning.

\subsection{Weight-Space Learning}

The idea of learning from neural network weights directly has gained traction with the availability of large model collections. Schürholt et al.~\cite{schurholt2022hyper} introduced hyper-representations---learned embeddings of model weights trained via contrastive learning. Their work demonstrated that weight-space features can predict properties like accuracy and robustness, establishing the feasibility of the paradigm. However, their architecture uses standard attention mechanisms without explicit permutation equivariance guarantees.

Unterthiner et al.~\cite{unterthiner2020predicting} pioneered predicting generalization from weights alone using simple statistics (mean, variance, norm of weight tensors). The Model Zoo dataset \cite{schurholt2022model} provided a benchmark with 50K+ models trained under varying conditions, enabling systematic evaluation of weight-space methods.

\subsection{Permutation-Equivariant Architectures}

Zhou et al.~\cite{zhou2023neural} introduced Neural Functional Networks (NFN), which process neural network weights using equivariant layers that respect the permutation symmetry of hidden neurons. Their NPLinear layer linearly transforms weight tensors while maintaining equivariance, and HNPPool produces permutation-invariant representations via pooling.

However, NFN's original evaluation focused primarily on generation rather than property prediction, and systematic comparison against matched-capacity non-equivariant baselines was limited. Our work addresses this gap by evaluating sample efficiency specifically.

\subsection{Inductive Biases and Symmetry}

Cohen and Welling~\cite{cohen2016group} showed that encoding symmetries via equivariant layers improves sample efficiency for image classification. Battaglia et al.~\cite{battaglia2018relational} argued that relational inductive biases are crucial for structured domains. Less understood is whether such biases are \emph{necessary} or merely \emph{convenient}. Our work provides empirical evidence that permutation invariance in weight-space learning cannot be learned from data regardless of scale.
